\section{Introduction}

\paragraph{Background.}
GPU kernel efficiency has become a central bottleneck in large-scale machine learning workloads~\cite{shah2024flashattention, zhang2025turbodiffusion, zhangefficient, zhangsageattention}. Prior work such as DeepSeek-V3~\cite{liu2024deepseek} demonstrates that at scale, competitive performance depends not only on model architecture but critically on kernel efficiency.
This has motivated a growing body of work using LLMs to automate Triton~\cite{tillet2019triton} kernel generation~\cite{yu2026towards}, ranging from training-based methods (AutoTriton~\cite{li2025autotriton}, TritonRL~\cite{woo2025tritonrl}) to agent-based iterative systems (GEAK~\cite{wang2025geak}, STARK~\cite{dong2025stark}) to search-and-reasoning approaches (KernelEvolve~\cite{liao2025kernelevolve}, ReGraph~\cite{gong2025large}).
\nocite{zhang2025spargeattn, zhang2025sla, zhang2026sla2, zhang2024sageattention2, zhang2025sageattention2++, zhang2025sageattention3, zhang2026sagebwd}

Accompanying this methodological progress, several benchmarks have been proposed. KernelBench~\cite{ouyang2025kernelbench} provides multi-level evaluation from operators to end-to-end models. TritonBench~\cite{li2025tritonbench} focuses on Triton kernels with a dual-channel testing framework. MultiKernelBench~\cite{wen2025multikernelbench} adds cross-platform evaluation, and Robust-KBench~\cite{lange2025towards} emphasizes robustness against misleading performance gains.

\paragraph{Limitations.}
Despite recent progress, two fundamental questions about LLM-based kernel generation remain unresolved.
First, the capability boundary is not well described: we do not yet know which types of tasks current methods handle reliably, which consistently fail, and why.
Second, the role of iterative refinement is not well understood: it is unclear whether different strategies improve compilation, correctness, or performance, and to what extent.
However, existing benchmarks are not designed to answer these questions, due to their unresolved task categories, insufficient correctness verification, and limited evaluation of efficiency.

\paragraph{Our Method.}
To address these limitations, we propose \textbf{\oursys}, built on TritonBench-T and extended in three directions: (1) a robust two-stage correctness protocol that rejects implementations passing output comparison by chance; (2) a unified 15-category taxonomy together with quantization and multi-precision task extensions, enabling fine-grained structural analysis; and (3) hardware-efficiency metrics beyond runtime.

\paragraph{Experimental Overview.}

Using \oursys, we conduct a systematic comparison of representative Triton kernel generation methods. Beyond aggregate benchmark results, we analyze method behavior across task categories, correctness outcomes, and efficiency metrics under a unified evaluation pipeline. We also collect error-correction and optimization pairs from the generation process for future training and inference-time improvement.

\paragraph{Contributions.}
Our main contributions are as follows:
\begin{itemize}
    \item We introduce \textbf{\oursys}, a benchmark for Triton kernel generation with category-aware evaluation of correctness and hardware efficiency.
    \item We conduct a systematic comparison of five representative methods under a unified pipeline.
    \item We identify three empirical findings that characterize the capability boundary of LLM-based kernel generation and provide mechanistic analysis for each.
    \item We release error--correction and optimization pairs collected during evaluation to support future training and inference.
\end{itemize}